% ============================================================
%  Paper 09: Strong Coupling alpha_s(M_Z) and Two-Loop Renormalization Group Running
%  Target: RevTeX 4-2 Compilation (SDGFT Series)
% ============================================================
\documentclass[amsmath,amssymb,aps,prd,twocolumn,superscriptaddress,nofootinbib]{revtex4-2}

\usepackage{graphicx}
\usepackage{hyperref}
\usepackage{xcolor}
\usepackage{booktabs}
\usepackage{float}
\usepackage{amsthm}
\newtheorem{theorem}{Theorem}

% ── SDGFT shared macros ──────────────────────────────────────
\usepackage{sdgft-macros}

% ── Document ─────────────────────────────────────────────────
\begin{document}

\preprint{SDGFT-2026-09}

\title{Strong Coupling alpha_s(M_Z) and Two-Loop Renormalization Group Running}

\author{David A. Besemer}
\email{david.besemer@sdgft.org}
\affiliation{Independent Researcher}

\date{\today}

\begin{abstract}
We derive the strong coupling constant alpha_s at the Z-boson scale from the topological face diagonal of the 24-cell lattice. We find the exact boundary condition alpha_s(M_Z) = sqrt(2)/12 ~ 0.117851. We implement a high-precision two-loop Runge-Kutta solver to verify the running of alpha_s up to the Planck scale, confirming asymptotic freedom.
\end{abstract}

\maketitle

\section{Introduction}
The strong coupling constant is the least precisely measured gauge coupling. We present a geometric derivation of its value at the Z mass.

\section{Theoretical Formulation}
Here we formulate the core mathematical relations governing the system:
\begin{equation}
  \als(M_Z) = \frac{\sqrt{2}}{12} \approx 0.117851
\end{equation}
\begin{equation}
  \mu \frac{d\als}{d\mu} = -\beta_0 \frac{\als^2}{2\pi} - \beta_1 \frac{\als^3}{8\pi^2} \quad (N_f = 5 \leftrightarrow 6 \text{ at } M_t)
\end{equation}
\section{Two-Loop Running and Thresholds}
We integrate the RGEs using a top-quark mass threshold Mt = 173.1 GeV. Below Mt, the number of active flavors is Nf=5, and above Mt, Nf=6.

\section{Conclusion}
The two-loop running shows stable convergence to the Planck scale, validating the strong coupling's geometric boundary condition.



\begin{thebibliography}{99}
\bibitem{Goldstein_Paper01} D. A. Besemer, ``Axiomatic Foundations of SDGFT,'' SDGFT-2026-01 (in preparation).
\bibitem{Goldstein_Paper04} D. A. Besemer, ``Effective Spectral Dimension and the Banach Fixed-Point Theorem in SDGFT,'' SDGFT-2026-04.
\bibitem{Goldstein_Code} D. A. Besemer, \texttt{sdgft} Python package, \url{https://github.com/david-besemer/sdgft} (2026).
\end{thebibliography}

\end{document}
